% !TeX root = ../../Book2.tex
% 纲要标题：### 9.2 外代数
% 纲要位置：计划纲要.md，第 1074 行。

\section{外代数}
\label{b2:sec:0902}

第八章的对称代数把生成元交换前后的差规定为零；外代数则把每个一次元素的平方规定为零，由此得到反对称关系。两者都是张量代数的商，却分别对应对称与交替的多线性运算。即使在特征二中，外代数的平方为零关系也不会消失。

\ProseParagraphBreak
不同次数的外幂不是彼此孤立的模。把输入首尾连接，可以在它们之间定义楔积，组成外代数。有基时，其计算归结为把指标排列成递增次序；没有基时，构造仍由交替映射的泛性质保证。本节把这两种描述联系起来，并核对它们在任意特征下都成立。

% 纲要要点（供目录核对）：
%
% * 基、维数与楔积；用交替坐标函数检测外幂基的系数
% * 分次交换律与奇数次平方为零；代数泛性质要求所有一次元素平方为零
% * 直和的外代数；分次张量积的符号与按次数分解
% * 域上纯楔积非零与线性无关的判据；一般环上非零与成基的区别
%

\subsection{基、维数与楔积}
\label{b2:subsec:0902:01}

设 \(\xi\in\bigwedge^pM\)、\(\eta\in\bigwedge^qM\)。希望用下式定义乘法：
\[
(m_1\wedge\cdots\wedge m_p)\wedge(n_1\wedge\cdots\wedge n_q)
=m_1\wedge\cdots\wedge m_p\wedge n_1\wedge\cdots\wedge n_q.
\]
纯楔积并非自由符号，因而需要先证明这个公式与它们之间的关系相容。

\begin{proposition}\label{b2:prop:wedge-product}
设 \(R\) 为交换环，\(M\) 为 \(R\) 模，\(p,q\ge0\)。规定两个纯楔积的积为
\[
(m_1\wedge\cdots\wedge m_p)\wedge(n_1\wedge\cdots\wedge n_q)
=m_1\wedge\cdots\wedge m_p\wedge n_1\wedge\cdots\wedge n_q.
\]
这一公式唯一延伸为 \(R\) 双线性映射
\[
\bigwedge^pM\times\bigwedge^qM\longrightarrow\bigwedge^{p+q}M,
\qquad(\xi,\eta)\longmapsto\xi\wedge\eta.
\]
这些映射满足结合律，\(R=\bigwedge^0M\) 中的 \(1\) 为单位元。因此
\[
\bigwedge M\coloneqq\bigoplus_{p\ge0}\bigwedge^pM
\]
成为分次 \(R\) 代数，称为 \(M\) 的\term[waidaishu]{外代数}[exterior algebra]；其乘法称为\term[xieji]{楔积}[wedge product]。
\end{proposition}
\begin{proof}
在 \(p=2,q=1\) 时，固定 \(n\in M\)，公式 \((m_1,m_2)\mapsto m_1\wedge m_2\wedge n\) 是交替双线性的。由\cref{b2:prop:exterior-alternating-universal}，它唯一地给出线性映射 \(L_n:\bigwedge^2M\rightarrow\bigwedge^3M\)，满足
\[
L_n(m_1\wedge m_2)=m_1\wedge m_2\wedge n.
\]
对 \(r,s\in R\) 及 \(n,n'\in M\)，\(L_{rn+sn'}=r L_n+s L_{n'}\) 在纯楔积上由多线性成立，故在整个 \(\bigwedge^2M\) 上成立。因此 \((\xi,n)\mapsto L_n(\xi)\) 双线性，且不依赖 \(\xi\) 的纯楔积表达。这就是把前两个变量先合成外幂，再对第三个变量使用一次外幂 \(\bigwedge^1M=M\) 的泛性质。

一般地，先设 \(p,q\ge1\)。固定 \(n_1,\ldots,n_q\)，右端关于 \(m_1,\ldots,m_p\) 交替且多线性，所以对 \(\xi\) 唯一线性延伸。所得公式关于各 \(n_j\) 仍多线性且交替：在纯 \(\xi\) 上已有这些等式，再由线性延伸到所有 \(\xi\)。第二次应用泛性质，即得关于 \(\eta\) 的唯一线性延伸。

若任一次数为零，则定义 \(r\wedge\xi=\xi\wedge r=r\cdot\xi\)，其中 \(r\in R=\bigwedge^0M\)。结合律在三个正次纯楔积上只是同一串因子的首尾连接，三线性与生成性使其在任意三个元素上成立；含零次因子时则由标量作用的结合律成立。\(1\) 因而为单位。直和中的元素只有有限多个非零次数分量，乘法扩张后的求和仍有限，故给出整个直和上的代数结构。
\end{proof}
\symidx{\(\bigwedge M\)：模 \(M\) 的外代数}

\begin{theorem}[有限自由模的外幂基]\label{b2:thm:exterior-power-basis}
设 \(R\) 为交换环，\(M\) 为有基 \(e_1,\ldots,e_n\) 的有限自由 \(R\) 模。对 \(0\le p\le n\)，全部递增指标楔积
\[
e_I\coloneqq e_{i_1}\wedge\cdots\wedge e_{i_p},
\qquad I=\{i_1<\cdots<i_p\}\subseteq\{1,\ldots,n\},
\]
组成 \(\bigwedge^pM\) 的基；\(p=0\) 时只有 \(e_{\varnothing}=1\)。对 \(p>n\)，有 \(\bigwedge^pM=0\)。特别地，在域上
\(\dim_k\bigwedge^p k^n=\binom np\)。
\end{theorem}
\begin{proof}
将各因子按 \(e_i\) 展开，再用多线性，任意纯楔积成为基向量楔积的有限线性组合。含重复指标的项为零，其余项由\eqref{b2:eq:alternating-permutation}化为带符号的递增指标项，故所列元素生成。当 \(p>n\) 时所有项都重复，得到零模。

为证明线性无关，希望对每个 \(e_I\) 构造一个只读取其系数的线性泛函，即在 \(e_I\) 上取一，在其他递增指标楔积上取零。先固定 \(1\le p\le n\) 及一个递增指标组 \(I\)。若 \(v_j=\sum_{a=1}^n x_{aj}e_a\)，每次从第 \(j\) 个输入选一个 \(I\) 中的基向量，按所选指标的排列符号求和，得到
\[
\Delta_I(v_1,\ldots,v_p)
=\sum_{\sigma\in S_p}\operatorname{sgn}(\sigma)
\prod_{j=1}^p x_{i_{\sigma(j)},j}.
\]
每个变量只出现一次，所以它多线性。若第 \(r,s\) 个输入相等，则把置换 \(\sigma\) 与 \(\sigma\circ(r\ s)\) 配对：二者不同，对应乘积相同而符号相反。每对的和为零，所以 \(\Delta_I\) 交替。这一成对消去在特征二中仍然有效。

泛性质给出线性映射 \(\widetilde\Delta_I:\bigwedge^pM\to R\)。对递增指标 \(J\)，当 \(I=J\) 时，求和中只有恒等置换对应的项为一，其余项均为零；当 \(I\ne J\) 时，每项都含有一个为零的坐标。因此 \(\widetilde\Delta_I(e_J)=1\) 当 \(I=J\)，否则为零。于是若 \(\sum_J a_Je_J=0\)，应用 \(\widetilde\Delta_I\) 便得 \(a_I=0\)。这证明线性无关，且没有除以任何整数。\(p=0\) 的结论由 \(\bigwedge^0M=R\) 直接得到。
\end{proof}

例如，在 \(R^3\) 中，二次外幂有基 \(e_1\wedge e_2,e_1\wedge e_3,e_2\wedge e_3\)。直接展开给出
\[
(e_1+2e_2)\wedge(e_2+e_3)
=e_1\wedge e_2+e_1\wedge e_3+2e_2\wedge e_3.
\]
系数 \(2\) 是否为零由所处的环决定，基的构造本身没有因此改变。

\subsection{分次交换律}
\label{b2:subsec:0902:02}

\begin{proposition}[分次交换律]\label{b2:prop:exterior-graded-commutativity}
设 \(R\) 为交换环，\(M\) 为 \(R\) 模，\(p,q\ge0\)。若 \(\xi\in\bigwedge^pM\)、\(\eta\in\bigwedge^qM\)，则
\begin{equation}\label{b2:eq:graded-wedge-sign}
\xi\wedge\eta=(-1)^{pq}\eta\wedge\xi.
\end{equation}
此外，每个奇数次齐次元素 \(\xi\) 都满足 \(\xi\wedge\xi=0\)，无须假设 \(2\) 可逆。
\end{proposition}
\begin{proof}
先取两个纯楔积。把前面的 \(p\) 个因子依次移过后面的 \(q\) 个因子，共交换 \(pq\) 次；每次由\cref{b2:prop:alternating-implies-skew}产生一个负号。这证明纯楔积情形，双线性再给出一般情形。

若 \(p\) 为奇数，分次交换律只直接给出 \(2\xi^2=0\)，在可能有二挠的环上不能据此推出平方为零。需要用外代数中重复因子为零的关系。写 \(\xi=\sum_{i=1}^t a_iw_i\)，其中 \(w_i\) 为 \(p\) 次纯楔积。每个 \(w_i\wedge w_i\) 都含重复因子，所以为零；不同指标的两项合为
\[
a_ia_j(w_i\wedge w_j+w_j\wedge w_i)=0
\]
，因为 \((-1)^{p^2}=-1\)。因此整个平方为零。
\end{proof}

若 \(2=0\) 于 \(R\)，分次交换律中的符号全部成为 \(1\)，外代数作为普通代数也是交换的，但仍保留每个一次元素的平方为零。例如
\[
\bigwedge(\mathbb F_2^2)
\cong\mathbb F_2[x,y]/(x^2,y^2),
\qquad e_1\longleftrightarrow x,\quad e_2\longleftrightarrow y.
\]
双方均由 \(1,x,y,xy\) 或相应楔积组成基，乘法关系一致，因而公式确为同构。它不同于无平方关系的对称代数 \(\mathbb F_2[x,y]\)。

外代数把每个一次元素的平方规定为零，而不只是把某一组基向量的平方规定为零。在线性映射的像中，\(u(x+y)^2=0\) 还要求两个交叉乘积之和为零。因此，即使 \(M\) 自由，仅检查基向量像的平方为零也不足以定义外代数上的同态。

\begin{proposition}[外代数的代数泛性质]\label{b2:prop:exterior-algebra-universal}
设 \(R\) 为交换环，\(M\) 为 \(R\) 模，\(A\) 为含单位元的结合 \(R\) 代数。若 \(R\) 线性映射 \(u:M\to A\) 满足 \(u(m)^2=0\) 对所有 \(m\in M\) 成立，则存在唯一 \(R\) 代数同态 \(\bigwedge M\to A\) 延伸 \(u\)。等价地，外代数是张量代数的商：
\[
\bigwedge M\cong T_R(M)/\langle m\otimes m\mid m\in M\rangle,
\]
分母表示这些二次元素生成的双边理想。
\end{proposition}
\begin{proof}
由假设，\(u(x)^2=u(y)^2=u(x+y)^2=0\)。由于 \(u\) 线性，展开第三个等式得
\[
0=(u(x)+u(y))^2
=u(x)^2+u(x)u(y)+u(y)u(x)+u(y)^2,
\]
因此 \(u(x)u(y)=-u(y)u(x)\)。乘积映射
\((m_1,\ldots,m_p)\mapsto u(m_1)\cdots u(m_p)\) 多线性；若两个输入相同，把其间的因子逐个交换，便出现一个 \(u(m)^2=0\)，所以它交替。逐次应用\cref{b2:prop:exterior-alternating-universal}得到各次外幂上的线性映射。首尾连接的公式保证它们合起来保持乘法和单位，故为代数同态。\(M\) 生成整个外代数，所以延伸唯一。

再由\cref{b2:thm:tensor-algebra-universal}，线性映射 \(M\to\bigwedge M\) 诱导满代数同态 \(T_R(M)\to\bigwedge M\)，并杀掉全部 \(m\otimes m\)，故经过上述商代数。反过来，商代数中的 \(M\) 的像都平方为零，由刚证明的泛性质得到反向代数同态。两复合在 \(M\) 与单位元上为恒等，因生成性而处处为恒等。
\end{proof}

\subsection{直和的外代数}
\label{b2:subsec:0902:03}

\begin{proposition}[直和的外幂]\label{b2:prop:exterior-direct-sum}
设 \(R\) 为交换环，\(M,N\) 为 \(R\) 模，\(p\ge0\)，则存在自然同构
\[
\bigoplus_{a+b=p}\left(\bigwedge^aM\otimes_R\bigwedge^bN\right)
\longrightarrow\bigwedge^p(M\oplus N),
\]
其在纯楔积上的公式为
\[
(m_1\wedge\cdots\wedge m_a)\otimes(n_1\wedge\cdots\wedge n_b)
\longmapsto
(m_1,0)\wedge\cdots\wedge(m_a,0)\wedge(0,n_1)\wedge\cdots\wedge(0,n_b).
\]
\end{proposition}
\begin{proof}
把 \(E=\bigwedge M\otimes_R\bigwedge N\) 按两侧次数之和分次。乘积原来的因子次序是 \(\xi,\eta,\xi',\eta'\)；要把 \(M\) 侧与 \(N\) 侧分别合在一起，就要让 \(N\) 侧次数为 \(b\) 的 \(\eta\) 越过 \(M\) 侧次数为 \(c\) 的 \(\xi'\)。在纯楔积上共发生 \(bc\) 次一次因子交换，所以应有符号 \((-1)^{bc}\)。据此采用带符号的分次张量积乘法，规定
\[
(\xi\otimes\eta)(\xi'\otimes\eta')
=(-1)^{bc}(\xi\wedge\xi')\otimes(\eta\wedge\eta'),
\qquad \deg\eta=b,\quad\deg\xi'=c.
\]
固定次数后此式分别线性且对两次张量关系平衡，故良定义。三因子结合律中，两种括号的符号指数分别为
\(bc+(b+d)e\) 与 \(de+b(c+e)\)，其中 \(d=\deg\eta'\)、\(e=\deg\xi''\)，二者相同；各外代数的结合律保证剩余因子也相同。单位是 \(1\otimes1\)。

映射 \(M\oplus N\to E\)、\((m,n)\mapsto m\otimes1+1\otimes n\) 的每个像都平方为零：两个同侧平方为零，两项交叉乘积分别为 \(m\otimes n\) 与 \(-m\otimes n\)。因此\cref{b2:prop:exterior-algebra-universal}给出分次代数同态 \(\alpha:\bigwedge(M\oplus N)\to E\)。

反方向，把 \(\xi\otimes\eta\) 送到两侧包含所诱导的楔积，得到线性映射 \(\beta:E\to\bigwedge(M\oplus N)\)。将中间的 \(\eta\) 与 \(\xi'\) 交换时，由\eqref{b2:eq:graded-wedge-sign}恰好得到 \((-1)^{bc}\)，所以 \(\beta\) 也是代数同态。两复合在所有 \(m\in M,n\in N\) 及单位元上都为恒等；这些元素生成双方代数，故 \(\alpha,\beta\) 互逆。取全次数 \(p\) 的分量，就得到命题中的直和同构。对 \(M,N\) 的线性映射，公式逐因子作用，故同构自然。
\end{proof}

一个 \(p\) 重楔积的输入来自 \(M\oplus N\)。将每个输入写成 \((m,n)=(m,0)+(0,n)\) 后按多线性展开，每一项恰有 \(a\) 个因子来自 \(M\)、\(b\) 个来自 \(N\)，其中 \(a+b=p\)。把 \(M\) 因子移到前面、\(N\) 因子移到后面，就落在上述直和的 \((a,b)\) 分量。因此直和公式是按两侧各取了多少个因子分类。特别地，
\[
\bigwedge^2(M\oplus N)
\cong\bigwedge^2M\oplus(M\otimes_RN)\oplus\bigwedge^2N.
\]
例如 \((m,0)\wedge(0,n)\) 属于中间项，而反向次序对应 \(-m\otimes n\)。这项负号已经由外代数的乘法决定，不是对直和分解额外指定的约定。若 \(M,N\) 分别具有 \(m,n\) 个基向量的有限自由模，比较基的数目便得到
\(\binom{m+n}{p}=\sum_{a+b=p}\binom ma\binom nb\)，其代数来源正是把一个递增指标组按两组基拆开。

\ProseParagraphBreak
在域上，纯楔积的非零性还能检测其因子是否独立。这里不需要先假定整个向量空间有限维：有限个因子只涉及一个有限维生成空间，再用直和公式比较它与母空间中的外幂。

\begin{proposition}\label{b2:prop:wedge-linear-independence}
设 \(k\) 为任意域，\(V\) 为 \(k\) 向量空间，\(p\ge1\)，\(v_1,\ldots,v_p\in V\)。则
\[
v_1\wedge\cdots\wedge v_p\ne0
\quad\Longleftrightarrow\quad
v_1,\ldots,v_p\text{ 线性无关}.
\]
\end{proposition}
\begin{proof}
若这些向量线性相关，取关系 \(\sum_{i=1}^p c_iv_i=0\) 中的非零系数 \(c_j\)，由域中可除性把 \(v_j\) 写成其余向量的线性组合。代入楔积并按多线性展开，每一项都含重复因子，因而为零。\(p=1\) 时这就是说 \(v_1=0\)，结论同样成立。

若这些向量线性无关，令 \(W=\operatorname{span}_k\{v_1,\ldots,v_p\}\)。它们构成有限维空间 \(W\) 的一组基，由\cref{b2:thm:exterior-power-basis}，\(v_1\wedge\cdots\wedge v_p\) 在 \(\bigwedge^pW\) 中是非零的基向量。由\cref{b2:thm:vector-basis-extension}将它们扩充为 \(V\) 的基，令 \(U\) 为其余基向量的线性生成空间，则 \(V=W\oplus U\)，即使 \(V\) 无限维也成立。\cref{b2:prop:exterior-direct-sum}给出
\[
\bigwedge^pV\cong
\bigoplus_{a+b=p}\bigl(\bigwedge^aW\otimes_k\bigwedge^bU\bigr).
\]
由该同构的纯楔积公式，\(W\hookrightarrow V\) 诱导的 \(\bigwedge^pW\rightarrow\bigwedge^pV\)，恰把 \(\bigwedge^pW\cong\bigwedge^pW\otimes_k k\) 送入 \((a,b)=(p,0)\) 的直和分量。因此它是单射，上述纯楔积在 \(\bigwedge^pV\) 中仍非零。
\end{proof}

在域上，独立向量还可以扩充成基；一般环上，非零楔积并不保证其因子能够扩充成母模的一组基。例如在 \(\mathbb Z^2\) 中，
\[
(2e_1)\wedge e_2=2(e_1\wedge e_2)\ne0,
\]
因为 \(e_1\wedge e_2\) 是自由 \(\mathbb Z\) 模 \(\bigwedge^2\mathbb Z^2\) 的基向量。但 \(2e_1,e_2\) 只能生成第一坐标为偶数的子模，不能生成整个 \(\mathbb Z^2\)。它们虽独立，楔积中的系数 \(2\) 却不是单位；“系数非零”与“系数可逆”的差别将在行列式的可逆性判据中再次出现。
